\part*{Part I. Compressed Analytic Wall Statement}
\addcontentsline{toc}{part}{Part I. Compressed Analytic Wall Statement}

\paragraph{Burden status at entry.}
Nothing primitive or carrier-specific has yet been certified. Part~I carries only the analytic-side burden of the paper: isolate the common vortex-stretching wall, show that every standard same-scope analytic route meets that wall, and compress the analytic search space into a finite atlas that Part~II can then reclassify. What remains after Part~I is \emph{not} a PDE continuation theorem, but the theorem-visible statement that any successful same-scope continuation argument would require a discriminator not supplied by the existing analytic bundle.

\paragraph{Compression note.}
The detailed route-by-route analytic atlas is moved to Appendix~E in this edition.
The main line retains only the shared bottleneck formulation and the atlas-closure theorem needed for
Part~II's pivot.

\section{The shared bottleneck (``the same wall'')}
We now formulate the common endpoint in a single statement.

\begin{proposition}[Common bottleneck formulation]\label{prop:common-bottleneck}
All standard analytic routes catalogued in Appendix~E reduce global regularity to establishing at least one of the following
\emph{critical upgrades} for every energy-class solution on $(0,T)$:
\begin{enumerate}[label=(\alph*)]
\item a scale-critical spacetime bound (Serrin--Prodi, $L^\infty_t L^3_x$, $L^5_{t,x}$, or equivalent Besov/Morrey control);
\item a scale-critical vorticity bound (e.g.\ $\int_0^T \|\omega\|_\infty<\infty$);
\item a rigidity/compactness principle excluding all nontrivial ancient solutions arising in Type II rescaling;
\item a depletion mechanism controlling $\mathcal T(u)$ without introducing extra hypotheses.
\end{enumerate}
In each case, the missing step is an \emph{unconditional} bridge from energy-class invariants to
critical control of the nonlocal vortex-stretching channel.
\end{proposition}

\section{Analytic obstruction and atlas closure}

\begin{theorem}[Analytic obstruction]
\label{thm:analytic-obstruction}
Within the classical formulation of the three--dimensional incompressible Navier--Stokes equations on $\mathbb T^3$, all known analytic continuation routes require an unconditional same--scope discriminator capable of globally suppressing vortex--stretching amplification.
No currently available analytic invariant, estimate, or continuation criterion provides such a discriminator.
\end{theorem}

\begin{proof}
Proposition~\ref{prop:common-bottleneck} isolates the shared endpoint of the standard routes: each reduces global regularity to an unconditional bridge from energy-class control to a scale-critical suppression of the vortex--stretching channel.
The route-by-route survey in Appendix~E shows that the standard criteria either assume critical control (smallness/integrability) or re-express the same missing bridge as an a priori blow-up condition.
None supplies an unconditional bundle-internal discriminator that closes the wall.
\end{proof}

\begin{corollary}[Analytic incompleteness]
\label{cor:analytic-incompleteness}
The classical analytic program cannot close the Navier--Stokes regularity problem without introducing a discriminator external to the existing analytic bundle.
\end{corollary}

\begin{definition}[Analytic continuation mechanism]\label{def:analytic-mechanism}
Fix the standard Navier--Stokes carrier on $\mathbb T^3$. An \emph{analytic continuation mechanism}
is a same-scope analytic attempt to certify continuation by strengthening control over some interior
feature of the admitted carrier while preserving its equations, scope constants, and licensed skins.
A purported mechanism that requires additional standing-bearing bundle data or a new same-scope
operator family is not an analytic continuation mechanism of the standard carrier but a scope change.
\end{definition}

\begin{theorem}[Atlas closure]
\label{thm:atlas-closure}
Within the classical Navier--Stokes carrier on $\mathbb T^3$, every analytic continuation mechanism in
the sense of Definition~\ref{def:analytic-mechanism} must modify exactly one of the following
structural features:
\begin{enumerate}[label=(\roman*), leftmargin=2.2em]
\item norm growth,
\item scale transfer,
\item geometric structure,
\item profile formation.
\end{enumerate}
No additional analytic continuation mechanism exists without introducing new bundle structure.
\end{theorem}

\begin{proof}
Energy/enstrophy programs target norm growth.
Littlewood--Paley and cascade analyses target scale transfer.
Geometric depletion and alignment arguments target geometric structure.
Blow-up profile and ancient-solution rigidity programs target profile formation.
A same-scope analytic continuation mechanism must act on some interior feature already visible in the
standard carrier skins. For that carrier, the compressed Part~I statement together with Appendix~E exposes only the four listed feature
families. Any purported fifth mechanism would require additional standing-bearing bundle data or a new
same-scope operator family not already licensed in the standard carrier. By definition, that is not a
continuation mechanism of the fixed carrier but a scope change. Hence the atlas is complete.
\end{proof}
\newpage
